\section{The certified first step: the hybrid theorem (package SHARP2-HYBRID)}\label{bp:p2s2hc}

Corollary~\ref{bp:p2s2:hybrid} (lever L6) states the hybrid principle with ``measured'' step-$0$ objects of the step identities; its numbers were
heuristic extrapolations (audits AUDIT-SHARP-V2-fixed-1 and AUDIT-SHARP-V2-iso-1, F6). This part gives the form a certificate can feed
(design document \texttt{docs/campaign/design/APOST-SHARP-V2-HYBRID.md}):
\begin{itemize}
\item \emph{Restart} (Proposition~\ref{bp:p2s2hc:restart}): the frozen conclusion for \emph{any} torus $K_1$ close to $K$ transfers to $K$. The hybrid
  applies a proved theorem to $K_1$; nothing inside the engine changes, and every theorem concluding the frozen body (v1, routes A, A$'$, B,
  the $\ell^1$ entry) can be restarted.
\item \emph{Transfer} (Lemmas~\ref{bp:p2s2hc:struct}--\ref{bp:p2s2hc:torsion}): every hypothesis of the v2 theorem for $K_1:=K+\Delta K$ follows from
  explicit bounds on $\Delta K$, on the linearised new error $\tilde E$, and on certified data of $K$.
\item \emph{The theorem} (Theorem~\ref{bp:p2s2hc:g1}) and \emph{the minimum with the existing routes} (Corollary~\ref{bp:p2s2hc:frozen}): every
  hypothesis is an explicit bound; the conclusion is the frozen body; never worse than the existing routes.
\end{itemize}
Soundness never uses that $\Delta K$ solves the cohomological equations. The producer computes $\Delta K$ (a trigonometric polynomial); the
checker computes norms of $\Delta K$ and the linearised new error. No small divisor is divided and $h$ is not composed with $K_1$ in the checker
for the step. The step is \emph{additive} (the form of route A, $\Delta K=L\xi^L+N\xi^N$ for the Newton correction): the modified step
$\tilde K\circ(\mathrm{id}+\xi^L)$ of route B is not needed, because the restart accepts any $K_1$, and would need a certified composition of models
(Remark~\ref{bp:p2s2hc:register}).

Norms are those of Part~\ref{bp:p2}: the max norm on vectors, the maximum absolute row sum on matrices; $T_r$ is the strip of width $r$.

%----------------------------------------------------------------------------------------------------
\subsection{Setting}\label{bp:p2s2hc:sec-set}

\begin{assumption}[Standing data of the hybrid]\label{bp:p2s2hc:standing}
\begin{enumerate}[label=(\alph*)]
\item $h$ and $\mathfrak B$ satisfy (H1) of Theorem~\ref{bp:p0:g1} with the bound $M$, and the refined data~\ref{bp:p2s2:R1} hold: $\abs{DX_h}\le M_1$,
  $\abs{D^2X_h}\le M_2$, $\abs{A\Omega_0-\Omega_0A}\le c_{Th}$ and $\abs{D(A\Omega_0-\Omega_0A)}\le c_{DTh}$ on $\mathfrak B$ ($A=DX_h$; defaults $2M_1$, $2M_2$).
\item $K$ satisfies (H2)--(H4) of Theorem~\ref{bp:p0:g1} on $T_{\rho_0}$ with the margin $\mu$ and the caps $\sigma_K,\sigma_B,\sigma_T$, and the refined
  data~\ref{bp:p2s2:R2}, \ref{bp:p2s2:R3} on $T_{\rho_0}$ with initial values $\abs{L^\top}\le\ell^\top$, $\abs N\le\nu$, $\abs{N^\top}\le\nu^\top$, and $\abs T\le t$
  on $T_{\rho_0-\delta_0}$. Write $\ell:=\sigma_K-\mu$, $\beta:=\sigma_B-\mu$, $m:=\sigma_T-\mu$ for the frozen initial bounds of $\abs{DK}$, $\abs B$, $\abs{\avg T^{-1}}$.
\item $\omega\in\Dio_{\gamma,\tau}$ with $0<\gamma\le\bar\gamma$.
\item A \emph{step} $\Delta K:\C^n\to\C^{2n}$, holomorphic on $T_{\rho_0}$, $1$-periodic ($\Delta K(\theta+k)=\Delta K(\theta)$ for $k\in\Z^n$) and real on $\R^n$.
  The \emph{new torus} is $K_1:=K+\Delta K$; its objects carry the index $1$ ($L_1=DK_1$, $B_1$, $N_1$, $T_1$, $E_1$).
\item A \emph{restart strip} $\rho_1\in(\rho_\infty,\rho_0]$ and a \emph{restart margin} $\mu_1\in(0,\mu)$; the restart bite $\delta_0':=\min\{1,(\rho_1-\rho_\infty)/6\}$
  and the torsion strip $T':=T_{\rho_1-\delta_0'}$ of the restarted run.
\end{enumerate}
\end{assumption}

\begin{definition}[Linearised new error]\label{bp:p2s2hc:lin}
With $E=X_h\circ K-DK\,\omega$ the invariance error of $K$,
\[
\tilde E(\theta):=E(\theta)+DX_h(K(\theta))\,\Delta K(\theta)-D\Delta K(\theta)\,\omega \qquad(\theta\in T_{\rho_0})
\]
(\Lean{KAM.Sharp2.Hybrid.linNewError}). It is affine in $\Delta K$ and is computed from the checker's models of $E$ and of $DX_h\circ K$ (both
produced by the composition of $h$ with $K$ that the checker already runs), the polynomial $\Delta K$ and its exact derivative. For the Newton
correction it is second order in $E$; a truncation of $\Delta K$ enters it linearly.
\end{definition}

\begin{definition}[Step bounds]\label{bp:p2s2hc:bounds}
Numbers $d_0,d_\infty,d_1,d_1^\top,d_0^\R,d_1^\R,d_1^{\R\top},\ell_E,\varepsilon_0',\varepsilon_0'^\R,o_0,o_0^\R,o_1,o_1^\R\ge0$ with
\begin{enumerate}[label=(S\arabic*)]
\item\label{bp:p2s2hc:S1} $\abs{\Delta K}\le d_0$ on $T_{\rho_1}$, and $d_0<\mu$;
\item\label{bp:p2s2hc:S2} $\abs{\Delta K}\le d_\infty$ on $T_{\rho_\infty}$;
\item\label{bp:p2s2hc:S3} $\abs{D\Delta K}\le d_1$ and $\abs{(D\Delta K)^\top}\le d_1^\top$ on $T_{\rho_1}$;
\item\label{bp:p2s2hc:S4} $\abs{\Delta K}\le d_0^\R$, $\abs{D\Delta K}\le d_1^\R$ and $\abs{(D\Delta K)^\top}\le d_1^{\R\top}$ on $\R^n$;
\item\label{bp:p2s2hc:S5} $\abs{\tilde E}\le\ell_E$ on $T_{\rho_1}$;
\item\label{bp:p2s2hc:S6} $\abs{DE}\le\varepsilon_0'$, $\abs{\Omega_K}\le o_0$ on $T'$ and $\abs{DE}\le\varepsilon_0'^\R$, $\abs{\Omega_K}\le o_0^\R$ on $\R^n$
  ($\Omega_K=L^\top\Omega_0L$, the frozen \Lean{lagrangeMat});
\item\label{bp:p2s2hc:S7} $\abs{\Omega_{K_1}}\le o_1$ on $T'$ and $\le o_1^\R$ on $\R^n$.
\end{enumerate}
Each bounds an explicit object: \ref{bp:p2s2hc:S1}--\ref{bp:p2s2hc:S4} a trigonometric polynomial and its derivative (weighted $\ell^1$ norms
of exact coefficients), \ref{bp:p2s2hc:S5} a model, \ref{bp:p2s2hc:S6} the derivative of the model of $E$ (Cauchy for its remainder,
Lemma~\ref{bp:p1:cauchy-strip}) and the exact polynomial $DK^\top\Omega_0DK$, \ref{bp:p2s2hc:S7} the exact polynomial $DK_1^\top\Omega_0DK_1$ (or
the defect lemma of the restarted run, Lemma~\ref{bp:p2s2:omega}, from $\abs{\eta^N_1}$). None is conclusion-shaped.
\end{definition}

%----------------------------------------------------------------------------------------------------
\subsection{The restart}\label{bp:p2s2hc:sec-restart}

\begin{proposition}[Restart]\label{bp:p2s2hc:restart}
Let $\mathbb B_\mu(K(\theta))\subset\mathfrak B$ for $\theta\in T_{\rho_\infty}$, let $K_1$ be any map with $\abs{K_1-K}\le d_\infty$ on $T_{\rho_\infty}$, and suppose the
conclusion of Theorem~\ref{bp:p0:g1} holds for $K_1$ with $(\rho_\infty,\mu_1,\sigma_{K,1})$ and the closeness $D_1$: there is $K_\infty$, holomorphic on
$T_{\rho_\infty}$, primary and real, with $E(K_\infty)=0$, $\mathbb B_{\mu_1/2}(K_\infty)\subset\mathfrak B$, $\abs{DK_\infty}\le\sigma_{K,1}-\mu_1/2$ and $\abs{K_\infty-K_1}\le D_1$ on
$T_{\rho_\infty}$. If
\[
d_\infty+D_1\le\mu/2\qquad\text{and}\qquad\sigma_{K,1}-\mu_1/2\le\sigma_K-\mu/2,
\]
then the conclusion of Theorem~\ref{bp:p0:g1} holds for $K$ with $(\rho_\infty,\mu,\sigma_K)$ and the closeness $d_\infty+D_1$, with the same $K_\infty$.
\source{Ours.}
\end{proposition}

\begin{proof}
\proofstep{Closeness and ball}\label{step:p2s2hc:rs-ball}
$\abs{K_\infty-K}\le\abs{K_\infty-K_1}+\abs{K_1-K}\le D_1+d_\infty\le\mu/2$ on $T_{\rho_\infty}$, so $\mathbb B_{\mu/2}(K_\infty(\theta))\subset\mathbb B_\mu(K(\theta))\subset\mathfrak B$.

\proofstep{The rest}\label{step:p2s2hc:rs-rest}
$\abs{DK_\infty}\le\sigma_{K,1}-\mu_1/2\le\sigma_K-\mu/2$; holomorphy, primarity, reality and invariance are properties of $K_\infty$ alone.
\end{proof}

No hypothesis of the theorem applied to $K_1$ is used: the proposition restarts any theorem that concludes the frozen body.

%----------------------------------------------------------------------------------------------------
\subsection{Transfer lemmas}\label{bp:p2s2hc:sec-transfer}

\begin{lemma}[The new torus]\label{bp:p2s2hc:struct}
Under Assumption~\ref{bp:p2s2hc:standing}, \ref{bp:p2s2hc:S1} and \ref{bp:p2s2hc:S3} with $d_0+\mu_1\le\mu$: $K_1$ is holomorphic and primary on $T_{\rho_1}$, real
on $\R^n$, $\mathbb B_{\mu_1}(K_1(\theta))\subset\mathfrak B$ on $T_{\rho_1}$, and $\abs{DK_1}\le\ell+d_1$ on $T_{\rho_1}$. With $\sigma_{K,1}:=\ell+d_1+\mu_1$,
$\abs{DK_1}\le\sigma_{K,1}-\mu_1$, and $\sigma_{K,1}-\mu_1/2\le\sigma_K-\mu/2$ if and only if $d_1+\mu_1/2\le\mu/2$.
\source{Ours.}
\end{lemma}

\begin{proof}
\proofstep{Structure}\label{step:p2s2hc:st}
Sums of holomorphic maps are holomorphic; $K_1(\theta+k)=K(\theta)+(k,0)+\Delta K(\theta)$; $K_1$ is real where $K$ and $\Delta K$ are;
$\mathbb B_{\mu_1}(K_1(\theta))\subset\mathbb B_{\mu_1+d_0}(K(\theta))\subset\mathbb B_\mu(K(\theta))$; $DK_1=DK+D\Delta K$ and $\abs{DK}\le\ell$.
\end{proof}

\begin{lemma}[The new error]\label{bp:p2s2hc:error}
Under Assumption~\ref{bp:p2s2hc:standing}, \ref{bp:p2s2hc:S1} and \ref{bp:p2s2hc:S5}: $E_1=X_h\circ K_1-DK_1\,\omega$ satisfies, on $T_{\rho_1}$,
\[
\abs{E_1}\le e_1:=\ell_E+\frac{M_2}2d_0^2 .
\]
\source{Lemma~\ref{bp:p1:taylor}; ours.}
\end{lemma}

\begin{proof}
\proofstep{The identity}\label{step:p2s2hc:er-id}
$E_1=\tilde E+\bigl(X_h(K+\Delta K)-X_h(K)-DX_h(K)\Delta K\bigr)$ pointwise, by Definition~\ref{bp:p2s2hc:lin}.

\proofstep{Taylor}\label{step:p2s2hc:er-tay}
The segment $[K(\theta),K_1(\theta)]$ lies in the convex ball $\mathbb B_\mu(K(\theta))\subset\mathfrak B$ since $d_0<\mu$; Lemma~\ref{bp:p1:taylor} with $A=M_2$ bounds the
bracket by $\frac{M_2}2\abs{\Delta K}^2$.
\end{proof}

\begin{lemma}[Frames of the new torus]\label{bp:p2s2hc:frames}
At a point $\theta$ where $G(\theta)$ is invertible, $\abs{L}\le\ell$, $\abs{L^\top}\le\ell^\top$, $\abs B\le\beta$, $\abs N\le\nu$, $\abs{N^\top}\le\nu^\top$, and $\Delta L:=D\Delta K(\theta)$
has $\abs{\Delta L}\le d_1$, $\abs{\Delta L^\top}\le d_1^\top$, put
\[
g:=\ell^\top d_1+d_1^\top\ell+d_1^\top d_1
\]
and assume $\beta g<1$. Then $G_1(\theta)=L_1^\top L_1$ is invertible and
\begin{gather*}
\abs{L_1^\top}\le\ell^\top+d_1^\top,\qquad\abs{B_1}\le\beta_1:=\frac{\beta}{1-\beta g},\qquad\abs{B_1-B}\le b:=\frac{\beta^2g}{1-\beta g},\\
\abs{N_1-N}\le n_\Delta:=\beta_1d_1+\ell b,\qquad\abs{N_1^\top-N^\top}\le n_\Delta^\top:=\beta_1d_1^\top+\ell^\top b .
\end{gather*}
\source{Lemma~\ref{bp:p2s:neumann}; ours.}
\end{lemma}

\begin{proof}
\proofstep{Gram increment}\label{step:p2s2hc:fr-g}
$G_1-G=L^\top\Delta L+\Delta L^\top L+\Delta L^\top\Delta L$, so $\abs{G_1-G}\le g$ (the row-sum norm is submultiplicative).

\proofstep{Neumann}\label{step:p2s2hc:fr-n}
$G_1=G(1+B(G_1-G))$ with $\abs{B(G_1-G)}\le\beta g<1$: Lemma~\ref{bp:p2s:neumann} gives the inverse and $\abs{B_1}\le\beta/(1-\beta g)$;
$B_1-B=-B_1(G_1-G)B$ gives $b$.

\proofstep{Normal frame}\label{step:p2s2hc:fr-nf}
$N_1-N=\Omega_0(\Delta L\,B_1+L(B_1-B))$ and $\abs{\Omega_0X}=\abs X$ ($\Omega_0$ is a signed permutation). $G$ and $G_1$ are symmetric, so $N^\top=BL^\top\Omega_0^\top$
and $N_1^\top-N^\top=(B_1\Delta L^\top+(B_1-B)L^\top)\Omega_0^\top$ with $\abs{X\Omega_0^\top}=\abs X$.
\end{proof}

\begin{lemma}[Torsion of the new torus]\label{bp:p2s2hc:torsion}
Let $U\subset T_{\rho_1}$, and assume the hypotheses of Lemma~\ref{bp:p2s2hc:frames} at every $\theta\in U$ with the same numbers; write $\nu_1:=\nu+n_\Delta$,
$\nu_1^\top:=\nu^\top+n_\Delta^\top$, $\ell_1:=\ell+d_1$, $\ell_1^\top:=\ell^\top+d_1^\top$. Let $\abs{DE}\le\varepsilon_0'$, $\abs{\Omega_K}\le o_0$, $\abs{DE_1}\le\varepsilon_1'$, $\abs{\Omega_{K_1}}\le o_1$
on $U$, and let $w_0,w_1$ bound the matrices $W=L^\top(A^\top+A)L-(DE)^\top L-L^\top DE$ of $K$ and of $K_1$ on $U$ (for instance
$w_j:=\ell_j^\top\abs{A^\top+A}\ell_j+\varepsilon_j'(2n\ell_j+\ell_j^\top)$ with $\abs{A^\top+A}\le4nM_1$). Then on $U$:
\begin{enumerate}[label=(\alph*)]
\item (frame part) $\abs{\mathcal T_{ab}(K_1)-\mathcal T_{ab}(K)}\le\Delta_{ab}:=n_\Delta^\top c_{Th}\nu_1+\nu^\top c_{DTh}\abs{\Delta K}\nu_1+\nu^\top c_{Th}n_\Delta$;
\item (local parts) $\abs{T-\mathcal T_{ab}(K)}\le\lambda_0:=\nu^\top\beta\varepsilon_0'+\beta^3w_0o_0$ and $\abs{T_1-\mathcal T_{ab}(K_1)}\le\lambda_1:=\nu_1^\top\beta_1\varepsilon_1'+\beta_1^3w_1o_1$;
\item (sup) $\abs{T_1}\le\min\{t+\Delta_{ab}+\lambda_0,\ \nu_1^\top c_{Th}\nu_1\}+\lambda_1$ where $\abs T\le t$;
\item (average) if $U\supset\R^n$ and $t_A\ge\sup_{\R^n}(\Delta_{ab}+\lambda_0+\lambda_1)$ with $mt_A<1$, then $\avg{T_1}$ is invertible and
  $\abs{\avg{T_1}^{-1}}\le m/(1-mt_A)$;
\item (Cauchy) under Lemma~\ref{bp:p2s2hc:error}, $\abs{DE_1}\le\varepsilon_1':=ne_1/\delta'$ on $T_{\rho_1-\delta'}$ for every $\delta'>0$.
\end{enumerate}
The torsion $T$ is the frozen one (Definition~\ref{bp:p2:objects}) at the frequency $\omega$; $\mathcal T_{ab}=BL^\top(A\Omega_0-\Omega_0A)LB$ is its frame part
(Lemma~\ref{bp:p2s:torsion-identity}, Lean \Lean{KAM.Sharp2.torsionAB}).
\source{Lemmas~\ref{bp:p2s:torsion-identity}, \ref{bp:p2s2:torsion-local}, \ref{bp:p2s2:torsion-ref}, \ref{bp:p1:cauchy-strip}; ours.}
\end{lemma}

\begin{proof}
\proofstep{Frame part}\label{step:p2s2hc:to-ab}
$BL^\top=N^\top\Omega_0$ and $LB=-\Omega_0N$, so $\mathcal T_{ab}=-N^\top\Omega_0\,C\,\Omega_0N$ with $C=A\Omega_0-\Omega_0A$ at $K(\theta)$. Telescope the triple product:
$X_1C_1Y_1-XCY=(X_1-X)C_1Y_1+X(C_1-C)Y_1+XC(Y_1-Y)$; $\abs{C_1-C}\le c_{DTh}\abs{\Delta K}$ by the increment bound along the segment in $\mathfrak B$.

\proofstep{Local parts}\label{step:p2s2hc:to-loc}
Lemma~\ref{bp:p2s2:torsion-local}'s estimates of $\mathcal T_F+\mathcal T_c$, at $K$ and at $K_1$ (Lean
\Lean{norm\_torsion\_sub\_torsionAB\_le} in \Lean{KAM.Sharp2}); the torsion
polynomial holds for $K_1$ on $T_{\rho_1}$ because $K_1$ is holomorphic there with values in $\mathfrak B$ and invertible Gram matrix
(Lemmas~\ref{bp:p2s2hc:struct}, \ref{bp:p2s2hc:frames}).

\proofstep{Sup and average}\label{step:p2s2hc:to-avg}
$T_1=T+(\mathcal T_{ab}(K_1)-\mathcal T_{ab}(K))-(T-\mathcal T_{ab}(K))+(T_1-\mathcal T_{ab}(K_1))$ gives the first bound in~(c), and
$T_1=\mathcal T_{ab}(K_1)+(T_1-\mathcal T_{ab}(K_1))$ with $\abs{\mathcal T_{ab}(K_1)}\le\nu_1^\top c_{Th}\nu_1$ the second. The average of $T_1-T$ over $[0,1]^n$ is bounded by its
sup on $\R^n$; Lemma~\ref{bp:p2s:neumann} as in step~\ref{step:p2s2:tr-n} (Lean \Lean{KAM.Sharp2.inv\_le\_of\_ref}).

\proofstep{Cauchy}\label{step:p2s2hc:to-c}
$E_1$ is holomorphic on $T_{\rho_1}$ and bounded by $e_1$ there; Lemma~\ref{bp:p1:cauchy-strip}.
\end{proof}

%----------------------------------------------------------------------------------------------------
\subsection{The hybrid theorem}\label{bp:p2s2hc:sec-g1}

\begin{definition}[Restart data and the hybrid condition]\label{bp:p2s2hc:data}
Given step bounds, Lemma~\ref{bp:p2s2hc:frames} with the strip data ($d_1,d_1^\top$; numbers $g,\beta_1,b,n_\Delta,n_\Delta^\top$) and with the real data
($d_1^\R,d_1^{\R\top}$; numbers with the superscript $\R$), the \emph{transferred initial values} of $K_1$ are
\[
\sigma_{K,1}:=\ell+d_1+\mu_1,\quad\sigma_{B,1}:=\beta_1+\mu_1,\quad\sigma_{T,1}:=\frac{m}{1-mt_A}+\mu_1,\quad
\ell_1^\top:=\ell^\top+d_1^\top,\quad\nu_1:=\nu+n_\Delta,\quad\nu_1^\top:=\nu^\top+n_\Delta^\top,
\]
$t_1:=$ the bound of Lemma~\ref{bp:p2s2hc:torsion}(c) on $T'$ (with $\varepsilon_1'=ne_1/\delta_0'$), $t_A:=\sup_{\R^n}(\Delta_{ab}^\R+\lambda_0^\R+\lambda_1^\R)$, and
$e_1:=\ell_E+\frac{M_2}2d_0^2$. The \emph{restart numbers} $N_1$ of a route $R\in\{\mathrm A,\mathrm A',\mathrm B\}$ are route $R$'s numbers
(Definition~\ref{bp:p2s2:data-ok}; Lean \Lean{NumA}, \Lean{NumB}) at the strip $\rho_1$ with the caps of the restarted run, its internal caps
and budgets (free, as for any run), and small-divisor constants~\ref{bp:p2s2:R4} at the bites of the schedule from $\rho_1$ (Lean
\Lean{NumA.restart}). The \emph{hybrid condition} $\mathrm{HybOk}_R$ is:
\begin{enumerate}[label=(\roman*)]
\item $d_0+\mu_1\le\mu$, $d_1+\mu_1/2\le\mu/2$, $\beta g<1$, $\beta g^\R<1$, $mt_A<1$;
\item route $R$'s numerical condition (Lean \Lean{routeAOk}, \Lean{routeAOkR} with the budgets \Lean{Ref.Budget}, \Lean{routeBOk}) for $N_1$ at some
  $X_\star^1>0$ with the initial values above, and its smallness at $e_1$ (routes A, A$'$: $e_1/(\gamma^4\delta_0'^{4\tau})<X_\star^1$; route B: $X_0(e_1)<X_\star^1$);
\item $d_\infty+D_1\le\mu/2$, with $D_1:=C_2^R(N_1,X_\star^1)\,e_1/(\gamma^2\rho_1^{2\tau})$ the closeness of route $R$'s branch at $N_1$.
\end{enumerate}
Every quantity is an explicit function of the step bounds, of certified data of $K$ and of route $R$'s numbers.
\end{definition}

\begin{theorem}[The hybrid theorem]\label{bp:p2s2hc:g1}
Assume Assumption~\ref{bp:p2s2hc:standing}, the step bounds~\ref{bp:p2s2hc:S1}--\ref{bp:p2s2hc:S7}, the small-divisor certificates~\ref{bp:p2s2:R4}
of $N_1$ on a set $\Omega\ni\omega$, and $\mathrm{HybOk}_R$. Then the conclusion of Theorem~\ref{bp:p0:g1} holds for $K$ with $(\rho_\infty,\mu,\sigma_K)$ and
the closeness $d_\infty+D_1$.
\source{Theorem~\ref{bp:p2s2:g1} (route $R$'s branch); ours.}
\end{theorem}

\begin{proof}
\proofstep{The hypotheses for $K_1$}\label{step:p2s2hc:g-hyp}
Lemma~\ref{bp:p2s2hc:struct}: (H2) for $K_1$ on $T_{\rho_1}$ with the margin $\mu_1$ and $\abs{DK_1}\le\sigma_{K,1}-\mu_1$. Lemma~\ref{bp:p2s2hc:frames} at every point of
$T_{\rho_1}\subset T_{\rho_0}$: (H3) with $\abs{B_1}\le\sigma_{B,1}-\mu_1$, and the refined data $\ell_1^\top,\nu_1,\nu_1^\top$. Lemma~\ref{bp:p2s2hc:torsion}(c) on $T'\subset T_{\rho_0-\delta_0}$
($\rho_1-\delta_0'\le\rho_0-\delta_0$ for $\rho_1\le\rho_0$): $\abs{T_1}\le t_1$; (d) on $\R^n$: (H4) with $\abs{\avg{T_1}^{-1}}\le\sigma_{T,1}-\mu_1$. Lemma~\ref{bp:p2s2hc:error}: the error
$e_1$ on $T_{\rho_1}$. (H1), $M_1$, $M_2$, $\omega$, $\gamma$ and the certificates are unchanged.

\proofstep{Restart}\label{step:p2s2hc:g-rs}
Route $R$'s branch at $N_1$ gives the frozen conclusion for $K_1$ with $(\rho_\infty,\mu_1,\sigma_{K,1})$ and the closeness $D_1$; Proposition~\ref{bp:p2s2hc:restart}
applies by $\mathrm{HybOk}_R$(i), (iii) and Lemma~\ref{bp:p2s2hc:struct}.
\end{proof}

\begin{corollary}[The minimum with the existing routes]\label{bp:p2s2hc:frozen}
Let $C_1^{\sharp\sharp},C_2^{\sharp\sharp}$ be the constants of Corollary~\ref{bp:p2s2:frozen-form}. For every $(\mathfrak B,h,K,\omega,\gamma,e)$ with the frozen hypotheses,
the refined data, the certificates, and a step with bounds~\ref{bp:p2s2hc:S1}--\ref{bp:p2s2hc:S7}: if
\[
C_1^{\sharp\sharp}e/(\gamma^4\rho^{4\tau})<1\qquad\text{or}\qquad\mathrm{HybOk}_R,
\]
then the conclusion of Theorem~\ref{bp:p0:g1} holds with the closeness $\max\{C_2^{\sharp\sharp}e/(\gamma^2\rho^{2\tau}),\,d_\infty+D_1\}$. It is never worse than the
existing routes: the first disjunct is Corollary~\ref{bp:p2s2:frozen-form}. Moreover $\Delta K=0$, $\rho_1=\rho_0$, $\mu_1=\mu$ make every increment vanish
and $e_1=e$, so $\mathrm{HybOk}_R$ is then route $R$'s own condition.
\source{Ours.}
\end{corollary}

\begin{proof}
\proofstep{By cases}\label{step:p2s2hc:fr}
The first case is Corollary~\ref{bp:p2s2:frozen-form}, the second Theorem~\ref{bp:p2s2hc:g1}; each closeness is at most the maximum.
\end{proof}

%----------------------------------------------------------------------------------------------------
\subsection{What the checker computes}\label{bp:p2s2hc:sec-checker}

\begin{remark}[Certifier interface]\label{bp:p2s2hc:checker}
Beyond the existing checks of $K$ (Hamiltonian bounds, embedding and ball, Gram witness, torsion and frame majorants, the invariance error
from the composition of $h$ with $K$), the checker:
\begin{enumerate}[label=(\arabic*)]
\item reads $\Delta K$ as a trigonometric polynomial (computed by the producer; any truncation is paid for in $\tilde E$);
\item forms the model of $\tilde E$: the model of $E$, plus the model of $DX_h\circ K$ (from the same composition) times $\Delta K$, minus the exact
  derivative of $\Delta K$ along $\omega$ (an interval for $\omega$); its weighted norm on $T_{\rho_1}$ is $\ell_E$ \ref{bp:p2s2hc:S5};
\item takes weighted $\ell^1$ norms of $\Delta K$ and $D\Delta K$ on $T_{\rho_1}$, $T_{\rho_\infty}$ and $\R^n$ \ref{bp:p2s2hc:S1}--\ref{bp:p2s2hc:S4};
\item bounds $DE$ on $T'$ and $\R^n$ (exact derivative of the polynomial part of $E$'s model, Cauchy for its remainder), and $\Omega_K$, $\Omega_{K_1}$
  exactly from the polynomials \ref{bp:p2s2hc:S6}, \ref{bp:p2s2hc:S7};
\item evaluates Definition~\ref{bp:p2s2hc:data} and route $R$'s numerical core for $N_1$ (the existing engines);
\item supplies small-divisor constants at $N_1$'s bites: for $\rho_1=\rho_0$ these are the plain certificate's; for $\rho_1=\rho_0-3\delta_0$ (and
  $\rho_0-\rho_\infty\le6$) the bites of its stages $\ge1$ and one more tail bite.
\end{enumerate}
No division by a small divisor and no composition with $K_1$ occur for the step.
\end{remark}

\begin{remark}[Register]\label{bp:p2s2hc:register}
\begin{enumerate}[label=(\roman*)]
\item \emph{Numbers} (\texttt{scripts/apost\_sharp/hybrid\_cert\_r1.py}; design document \S4). They are reported in the design document; they
  are design numbers in the enclosures' norms, computed in double precision, not certificates.
\item \emph{Relation to Corollary~\ref{bp:p2s2:hybrid}.} Superseded: the restart applies the whole theorem to $K_1$, so neither the step identities
  nor stage-$0$ constants are needed, and the step need not be a Newton step. The engine of Proposition~\ref{bp:p2s2:engine} is not modified.
\item \emph{Restart schedule.} The restarted run recomputes its bite from $\rho_1$. For $\rho_1=\rho_0-3\delta_0$ and $\rho_0-\rho_\infty\le6$ this is the original
  schedule from stage $1$ (cf.\ AUDIT-B10D-1, F5); for $\rho_1=\rho_0$ it is the original schedule itself.
\item \emph{The step's form.} A route-B step $(K+N\xi^N)\circ(\mathrm{id}+\xi^L)$ would need a certified composition of models with a shift of the argument
  and its Taylor remainder in $\theta$; the measured new error of the additive step is already second order (design document \S4), so it is
  not used.
\item \emph{Observation on R1.} In R1's rotated frame $H$ is linear in $P_2$, and the frozen torsion average is singular: $\avg T=\mathrm{diag}(1.96,0)$
  (sampled), while the bordered torsion is invertible with inverse $0.584$. Only the isoenergetic theorem applies to R1; its fixed-frequency
  rows in the design tables are formal. The hybrid for R1 is the isoenergetic one (track~I part \texttt{p2s2-iso-hybrid-cert}).
\item \emph{Suspected errata.} None in a pinned source. \emph{Deviations.} None from a source; the restart and the transfer lemmas are ours.
\end{enumerate}
\end{remark}
